\documentclass{article}
\usepackage[utf8]{inputenc}
\usepackage[spanish]{babel}
\usepackage{amsmath}
\usepackage{graphicx}
\usepackage{circuitikz}
\usepackage{subcaption}
\usepackage{geometry}
\usepackage{enumitem}
\usepackage{comment}
\geometry{a4paper, margin=1in}

\title{\textbf{PRÁCTICA 26 \\ TEOREMA DE THÉVENIN}}
\author{}
\date{}

\begin{document}

\maketitle
\thispagestyle{empty}

\begin{comment}
\begin{center}
\textbf{154 Prácticas de electricidad}
\end{center}
\end{comment}

% \section*{Introducción}

\begin{figure}[h!]
    \centering
    \begin{subfigure}[b]{0.45\textwidth}
        \centering
        \begin{circuitikz}[scale=0.8]
            \draw (0,0) to[V, l=60V] (0,4)
                  (0,4) to[R, l=$40\Omega$] (3,4)
                  (3,4) to[R, l=$60\Omega$] (6,4)
                  (6,4) to[short] (6,0)
                  (6,0) to[short] (0,0)
                  (3,4) to[R, l=$120\Omega$] (3,0)
                  (3,0) to[R, l=$160\Omega$] (6,0)
                  (3,4) to[R, l=$50\Omega$, *-*] (3,0);
            \node at (3,4.3) {B};
            \node at (3,-0.3) {C};
        \end{circuitikz}
        \caption{}
    \end{subfigure}
    \hfill
    \begin{subfigure}[b]{0.45\textwidth}
        \centering
        \begin{circuitikz}[scale=0.8]
            \draw (0,0) to[V, l={$V_{TH}=8V$}] (0,2)
                  (0,2) to[R, l={$R_{TH}=72\Omega$}] (0,4)
                  (0,4) to[short, -o] (2,4) node[right] {B}
                  (0,0) to[short, -o] (2,0) node[right] {C}
                  (2,3.5) to[R, l={$R_5=100\Omega$}] (2,0.5);
        \end{circuitikz}
        \caption{}
    \end{subfigure}
    
    \vspace{0.5cm}
    
    \begin{subfigure}[b]{0.45\textwidth}
        \centering
        \begin{circuitikz}[scale=0.8]
            \draw (0,0) to[V, l=60V] (0,4)
                  (0,4) to[R, l=$40\Omega$] (3,4)
                  (3,4) to[R, l=$60\Omega$] (6,4)
                  (6,4) to[short] (6,0)
                  (6,0) to[short] (0,0)
                  (3,4) to[R, l=$120\Omega$] (3,0)
                  (3,0) to[R, l=$160\Omega$] (6,0)
                  (3,4) to[open, v=$V_{TH}$] (3,0);
            \node at (3,4.3) {B};
            \node at (3,-0.3) {C};
        \end{circuitikz}
        \caption{}
    \end{subfigure}
    \hfill
    \begin{subfigure}[b]{0.45\textwidth}
        \centering
        \begin{circuitikz}[scale=0.8]
            \draw (0,4) to[short] (0,3)
                  (0,3) to[R, l=$40\Omega$] (0,1)
                  (0,1) to[R, l=$160\Omega$] (0,0)
                  (0,0) to[short] (3,0)
                  (3,0) to[R, l=$120\Omega$] (3,2)
                  (3,2) to[R, l=$60\Omega$] (3,4)
                  (3,4) to[short] (0,4)
                  (0,4) to[short, -o] (4,4) node[right] {B}
                  (0,0) to[short, -o] (4,0) node[right] {C};
        \end{circuitikz}
        \caption{}
    \end{subfigure}
    
    \caption{Fig. 26-2 — Circuito puente desequilibrado.}
    \label{fig:unbalanced_bridge}
\end{figure}

\section*{Verificación experimental del teorema de Thévenin}
Es posible determinar por medición los valores de \( V_{TH} \) y \( R_{TH} \) para una carga \( R_L \) en una red dada. Luego podemos ajustar experimentalmente la salida de una fuente de alimentación de tensión regulada en \( V_{TH} \) conectando una resistencia cuyo valor sea \( R_{TH} \), en serie con \( V_{TH} \) y con \( R_L \). Podemos medir \( I \) en este circuito equivalente. Si el valor medido de \( I_L \) en \( R_L \) de la red original es el mismo que el de \( I \) medido en el equivalente de Thévenin, tenemos una comprobación del teorema de Thévenin. Para una comprobación más completa este proceso tendría que ser repetido varias veces con circuitos cualesquiera.

\section*{RESUMEN}
\begin{enumerate}[label=\arabic*.]
    \item El teorema de Thévenin enuncia que cualquier red de dos terminales puede ser reemplazada por un circuito sencillo equivalente que actúa como el circuito original en la carga conectada a los dos terminales.
    
    \item El circuito equivalente se compone de un generador o fuente de alimentación llamado generador de Thévenin (\( V_{TH} \)) en serie con la resistencia interna de Thévenin (\( R_{TH} \)), en serie a su vez con la carga, entre los dos terminales en cuestión. Así para el circuito complejo de la figura 26-2 a, el equivalente de Thévenin en los terminales BC es el de la figura 26-2 b.
    
    \item Para determinar la tensión de Thévenin \( V_{TH} \) se abre la carga en los dos terminales en la red original y se calcula la tensión en estos dos terminales. Esta tensión con carga abierta es \( V_{TH} \).
    
    \item Para determinar la resistencia de Thévenin \( R_{TH} \) se mantiene abierta la carga en los dos terminales en la red original y se reemplaza la fuente de alimentación original con su propia resistencia interna. Luego se calcula la resistencia entre los terminales de carga abierta, hacia el circuito original.
    
    \item El teorema de Thévenin es aplicable a circuitos lineales con fuentes independientes y dependientes, y es especialmente útil para analizar circuitos con diferentes valores de carga.
\end{enumerate}

\end{document}